\documentclass[11pt]{amsart}

\usepackage[T1]{fontenc}
\usepackage[utf8]{inputenc}
\usepackage{microtype}
\usepackage{mathtools}
\usepackage{amssymb}
\usepackage{enumitem}
\usepackage[colorlinks=true,citecolor=blue,linkcolor=blue,urlcolor=blue]{hyperref}

\newtheorem{theorem}{Theorem}[section]
\newtheorem{proposition}[theorem]{Proposition}
\newtheorem{lemma}[theorem]{Lemma}
\newtheorem{corollary}[theorem]{Corollary}
\theoremstyle{definition}
\newtheorem{remark}[theorem]{Remark}

\newcommand{\CharVar}{\operatorname{Char}}

\date{\today}

\begin{document}
\section{Simple add-on why cyclicity is an equivalent statement}
As is already written correctly in the main file, Stafford showed that the Conjecture 3.8 of \cite[p.~438]{Stafford1978}
is equivalent to stable cyclicity of every finitely generated torsion right $A_n$-module \cite[pp.~437--438]{Stafford1978}. Here, I tried to provide this argument in a self-contained way that is understandable for non-experts (such as myself). However, I just managed to understand the $"\Rightarrow"$ direction, which is anyway the one that is interesting. 
Statement in question:\\
\begin{lemma}[Equivalence to stably cyclic torsion modules]
The following statements are equivalent:
\begin{enumerate}
    \item For every $d\neq 0 \in A_n$, there exists an $f \in A_n$, such that $A_n = dA_n + fdA_n$
    \item Every finitely generated torsion right $A_n$-module is stably cyclic. 
\end{enumerate}
\end{lemma}
\begin{proof}
   "$\Rightarrow$": \\
        Suppose for every $d\neq 0 \in A_n$, there exists an $f \in A_n$, such that $A_n = dA_n + fdA_n$. Let us denote a torsion module by $T$. Stafford's theorem (Theorem 3.7 \cite[p.~438]{Stafford1978}) established, that every finitely generated torsion module can be generated by 2 generators, so let us denote the generators of $T$ with $t_1$ and $t_2$. We can define a homomorphism 
        \begin{equation}
            \phi:A_n \oplus A_n\rightarrow T \quad \phi(r,s)\mapsto t_1r+t_2s.
        \end{equation}
        This isomorphism is surjective, since $t_1, t_2$ generate all of $T$. By the first isomorphism theorem (for $\phi:M\rightarrow N$, $im(\phi) \cong M/ker(\phi)$), it follows 
        \[T \cong (A_n\oplus A_n)/M\]
        with $M = ker(\phi)$. \\
        $T$ is torsion and therefore $\forall t\in T: \exists m \in A_n\backslash\{0\}: tm = 0.$ Specifically, let us denote with $d_1$ and $d_2$ the annihilators of $t_1$ and $t_2$, i.e.
        \[t_1d_1 = t_2d_2 = 0.\]
        $A_n$ is a right Ore domain, that means it has the property that for any nonzero $d_1, d_2$ there exists elements $r_1, r_2 \in A_n$ such that 
        \[d_1r_1 = d_2r_2 \neq 0.\]
        Defining therefore $d := d_1r_1 = d_2r_2$ yields us a common annihilator of $t_1$ and $t_2$.\\
        $(d, 0)$ and $(0,d)$ are elements in $M = ker(\phi)$, since $\phi(d, 0) = t_1d = 0$ and $\phi(0, d) = t_2d = 0.$\\
        From $A_n = dA_n + fdA_n$ it follows that there exist $r,s \in A_n$ such that $1 = dr + fds$. Because $(0,d), (d, 0) \in M$, also 
        $(d, 0)s + (0, d)r = (ds, dr) = (ds, 1-fds) =: (a, 1-fa) \in M$, i.e. $(a, 1-fa)$ is an element of $M$. \\
        On the other hand, $A_n\oplus A_n$ can be generated by $(a, 1-fa)$ and $(1, -f)$. We can check this by finding $u,v$ such that for all $(x,y)\in A_n\oplus A_n$ we can write 
        \[(x,y) = (a, 1-fa)u + (1, -f)v.\]
        Suitable choices are $u = y+fx$ and $v = x-ar$.\\
        Lastly, since  $A_n\oplus A_n$ can be generated by $(a, 1-fa)$ and $(1, -f)$ and $(a, 1-fa)$ is an element of $M$, it follows that $T \cong /A_n\oplus A_n)/M$ can be generated by $(1, -f)$, i.e. $T$ is cyclic. 
        %% left: stably cyclic
\end{proof}

\end{document}